\chapter{Ontological Engineering}

\chapterepigraph{%
  An ontology is not a description of reality.\\
  It is a bet on which distinctions matter.\\
  Good ontologies make accurate bets.\\
  Bad ontologies lose them quietly.
}{Flyxion, \textit{Semantic Infrastructure}}

\sidenote{Connects to\\semantic modules\\(Ch.\,25) and\\obstruction theory\\(App.\,I)}

Classification (Chapter~45) produces taxonomies. Ontological engineering
goes further: it designs the category systems themselves, making explicit
choices about which distinctions to encode, which to elide, and how
categories relate to each other. This chapter develops the theory of
ontological engineering using the semantic module framework.

\section{Building Categories}

Every ontology is a system of semantic modules (Chapter~25): each category
is a module with a content (what entities belong to it), an admissibility
structure (what properties make something a member), and an interface
(how the category relates to other categories).

\begin{definition}{Ontological System}{ontological-system}
  An \emph{ontological system} is a tuple $(\mathcal{C}, \mathcal{R},
  \mathcal{A})$ where:
  \begin{itemize}
    \item $\mathcal{C} = \{C_1, \ldots, C_n\}$ is a set of categories
          (semantic modules),
    \item $\mathcal{R}$ is a set of relations between categories
          (is-a, part-of, caused-by, associated-with),
    \item $\mathcal{A}$ is the joint admissibility structure
          (which combinations of category membership are consistent).
  \end{itemize}
\end{definition}

\section{Revising Categories}

Ontologies must be revised as knowledge evolves. Revision is not merely
adding new categories — it may require restructuring existing ones,
splitting overloaded categories, merging categories that were wrongly
distinguished, or changing the underlying relations.

\begin{definition}{Ontological Revision}{onto-revision}
  An \emph{ontological revision} is a morphism $r : \mathcal{O}_1 \to
  \mathcal{O}_2$ between ontological systems that: (1) is admissibility-
  preserving (valid combinations in $\mathcal{O}_1$ remain valid in
  $\mathcal{O}_2$), and (2) minimizes the disruption to existing
  knowledge (minimal semantic distance between corresponding categories).
\end{definition}

The Copernican revision of astronomy is an ontological revision:
the categories ``fixed Earth'' and ``moving Sun'' were replaced by
``orbiting Earth'' and ``central Sun.'' The revision was valid
(observational data remained consistent) and disruptive (the category
structure changed fundamentally).

\section{When Categories Fail and How to Fix Them}

Ontological systems fail in predictable ways:

\begin{itemize}
  \item \textbf{Over-splitting:} two categories that should be one
        (``heat'' and ``caloric'' before thermodynamics unified them).
  \item \textbf{Under-splitting:} one category that should be two
        (``atom'' before the discovery of the nucleus).
  \item \textbf{Wrong relations:} correct categories, wrong relationships
        (``evolution'' before the mechanism of natural selection).
  \item \textbf{Missing categories:} the ontology lacks a category for
        an important class of phenomena (``quantum field'' before
        quantum field theory).
\end{itemize}

Each failure type has a signature in the accessibility landscape:
over-splitting creates unnecessary wells; under-splitting hides
peaks; wrong relations produce incorrect gradient flows; missing
categories leave entire regions of the landscape unnamed and inaccessible.

\begin{namedthm}{The Ontological Engineering Principle}
  A good ontology minimizes total constraint violation across the
  knowledge infrastructure it serves: it encodes the distinctions
  that matter (reducing $\kappa$ for accurate prediction), ignores
  the distinctions that don't (reducing maintenance cost), and
  provides relations that support productive inference (maximizing
  $S_{\mathbf{Sem}}$).
\end{namedthm}

\begin{exercise}
  The ICD (International Classification of Diseases) is a medical
  ontology used globally for clinical and epidemiological purposes.
  Identify one case of over-splitting, one of under-splitting,
  and one of wrong relations in the current version. For each,
  propose a revision and analyze its accessibility consequences.
\end{exercise}

\begin{exercise}
  Design a minimal ontology for the concept of ``intelligence.''
  Your ontology should: (1) encompass biological and artificial
  intelligence, (2) encode the key distinctions (general vs. narrow,
  adaptive vs. fixed), and (3) be revision-friendly. What categories
  and relations does it include? What potential failure modes does it have?
\end{exercise}
